\chapter{Static Single Assignment}
\label{ch:ch12-static-single-assignment}

Every dataflow fact so far has a name problem underneath it.
A variable assigned in five blocks carries five histories under one name, and each use must sift which histories reach it.
Reaching definitions answers by sets, liveness answers by backward sets, and clients pay the set cost at every query.
The problem of this chapter is to rename variables so that each name has exactly one definition, and to make merging explicit where paths join.
After renaming, a use refers to one definition by construction, and many flow queries become local lookups.
Section~\ref{sec:ch12-meaning} defines the form, Section~\ref{sec:ch12-placement} places the merges, and Section~\ref{sec:ch12-worked} renames the graph in Figure~\ref{fig:ch12-ssa} by hand.

\section{What single assignment means}
\label{sec:ch12-meaning}

A program is in \term{static single assignment} form when every variable has exactly one static definition and every use refers to that definition.
Straight-line code reaches the form by versioning: \code{x = 1} becomes \code{x1 = 1}, \code{x = x + 2} becomes \code{x2 = x1 + 2}, and later uses take the current version.
Branches complicate versioning because two arms define different versions of the same source variable.
Joins therefore need a merge operation that selects the version from the arm actually taken.
That operation is the \term{phi function}, written $\assignphi{x_3}{x_1,x_2}$ to mean that $x_3$ takes $x_1$ from the left edge and $x_2$ from the right edge.
Phi functions sit at the top of join blocks, execute conceptually on edge traversal, and disappear before machine code emission.
The \enquote{static} in the name marks that uniqueness holds at compile time over program text, not at run time over executed stores.

\begin{irlisting}
; Source before SSA: x defined on both arms, used at join
; After SSA: versions x1 x2, merge x3 by phi at B4 head
B1: br c, B2, B3
B2: x1 = 1
B2: jmp B4
B3: x2 = 2
B3: jmp B4
B4: x3 = phi(x1, x2)
B4: y1 = x3 + 1
B4: use y1
\end{irlisting}

The listing shows the three ingredients together: versioned definitions, one phi per merged variable at the join, and uses rewritten to current versions.
The phi has one argument per incoming edge, ordered by edge, not by value.
A phi at a three-way join needs three arguments.
A loop-carried variable needs a header phi when its incoming definitions must be merged; a variable overwritten before every use need not have one in pruned form.
Variables never assigned keep no versions and need no phis.
Memory locations need no phis until the compiler promotes them to registers, because stores through pointers do not create unambiguous definitions.
The renaming in Section~\ref{sec:ch12-implementation} turns this picture into an algorithm, and Table~\ref{tab:ch12-df} supplies the placement facts it needs.

Minimality matters because extra phis cost time and space without adding precision.
A program is in \term{minimal SSA} form when it holds exactly the phis needed to preserve single-assignment meaning, no more.
A program is in \term{pruned SSA} form when it additionally drops phis for variables never live across the join.
Pruned form is what compilers build, since dead phis would immediately burden the next passes.
Semi-pruned construction filters out variables that are never live on entry to any block, avoiding a full per-block liveness filter while retaining some dead phis.
The placement algorithm below targets minimal form, and a liveness filter trims it to pruned form afterward.
Appel develops the renaming stack that makes this pipeline concrete, see~\cite{appel-1998-modern}.

\section{Where phis must go}
\label{sec:ch12-placement}

A phi for variable $v$ belongs at every join where two distinct definitions of $v$ arrive along different paths.
Dominance frontiers compute exactly those joins from the dominator tree.
A node $n$ \term{dominates} $m$ when every path from entry to $m$ passes through $n$.
The \term{dominance frontier} $\DF(n)$ contains each block $y$ for which $n$ dominates a predecessor of $y$ but does not strictly dominate $y$.
If $n$ defines $v$ and $j$ sits in $\DF(n)$, then some path carries that definition to $j$ alongside a path that avoids $n$, so $j$ needs a phi for $v$.
Iterating this rule to closure over newly added phis yields the minimal placement.
The iteration stops because each block gains at most one phi per variable.
A \term{dominator tree} records immediate dominators as parent links, with the entry as root and every block hanging below its strictest dominator.
Frontiers fall out of the tree by one upward walk per edge: for edge $a$ to $b$, climb from $a$ while the current node does not strictly dominate $b$, adding $b$ to each visited frontier.
Different edges can repeat parts of an upward walk, so this construction does not have a general linear-time bound.
Loop heads appear in their own frontiers when the back edge climbs past them, which is why loop variables need phis even without any visible branch.
Figure~\ref{fig:ch12-ssa} shows a loop head with exactly that self-frontier, and Table~\ref{tab:ch12-df} lists the frontier sets behind the placement.
Morgan treats dominator-based placement as the standard front end for SSA construction, see~\cite{morgan-1998-optimizing}.

\begin{algorithm}
\caption{Minimal phi placement by iterated dominance frontiers}\label{alg:ch12-placement}
\begin{pseudocode}
\Require CFG $G$, dominator tree with $\DF(n)$ per node, set $\Def(v)$ of blocks defining each $v$
\Ensure one phi for $v$ at each join in the iterated frontier of $\Def(v)$
\ForAll{variable $v$ defined anywhere in $G$}
\State $Work \gets \Def(v)$
\State $Placed \gets \emptyset$
\While{$Work \neq \emptyset$}
\State remove some $n$ from $Work$
\ForAll{$j \in \DF(n)$}
\If{$j \notin Placed$ for $v$}
\State insert $\assignphi{v_j}{...}$ at the head of $j$
\State $Placed \gets Placed \cup \{j\}$
\If{$j \notin \Def(v)$ originally}
\State add $j$ to $Work$
\EndIf
\EndIf
\EndFor
\EndWhile
\EndFor
\end{pseudocode}
\end{algorithm}

Algorithm~\ref{alg:ch12-placement} processes one variable at a time over definition sites and frontier edges.
Each definition site seeds the worklist, each frontier join gains a phi, and each fresh phi acts as a definition that can force further phis downstream.
The test against $Placed$ prevents duplicate phis at the same join for the same variable.
The test against original definitions prevents reprocessing blocks that already seeded the search.
The result is the conventional unpruned iterated-frontier placement; it can include phis whose results are dead.
Pruned placement additionally requires the variable to be live on entry to the candidate block, or removes dead phi chains afterward.
Treat phi arguments as uses on their corresponding incoming edges, not simultaneous uses at the start of the join block.
An implementation should place phis for all variables before any renaming, because renaming while placing mixes two traversals with different invariants.

Renaming walks the dominator tree with one version stack per source variable.
It rewrites each definition to push a fresh version, each use to read the stack top, and each phi argument in a successor to read the stack top along that edge.
It recurses into tree children with the pushed versions visible, then pops them on return so siblings see the correct prefix.
Seed parameter and other entry definitions before traversal.
An empty version stack at a reachable use signals a missing definition or a language-specific uninitialized value; inventing an arbitrary version is not a sound repair.
The walk is linear in program size plus phi arguments, with stack depth bounded by tree height.
A subtle bug rewrites phi arguments as if they executed at the join head rather than on incoming edges.
The fix reads the predecessor-specific stack state per edge, which the tree walk holds by processing the predecessor before descending.
Section~\ref{sec:ch12-worked} traces the stacks on a small graph so the edge rule becomes concrete.
\section{Worked example with loop and branch}
\label{sec:ch12-worked}

We place and rename by hand on Figure~\ref{fig:ch12-ssa}, which holds a branch inside a loop.
Variable \code{x} is defined in $B1$ before the loop, in $B3$ on the left arm, and in $B4$ on the right arm.
For a separate pruning experiment, suppose $B3$ also defines an unused variable \code{y}.
The dominator tree has $B1$ as root, $B2$ as its child, $B3$, $B4$, and $B5$ as children of $B2$, and $B6$ as a child of $B5$.
The branch arms have frontier $\{B5\}$; $B2$ and $B5$ have frontier $\{B2\}$.
The entry has an empty frontier.
Table~\ref{tab:ch12-df} lists these sets exactly as the placement loop consumes them.

\begin{figure}
\begin{tikzpicture}
\node[basic block] (b1) {B1: x = 0};
\node[basic block, below=1.0cm of b1] (b2) {B2: head phi x};
\node[basic block, below left=1.0cm and 0.6cm of b2] (b3) {B3: x = x + 1};
\node[basic block, below right=1.0cm and 0.6cm of b2] (b4) {B4: x = x + 2};
\node[basic block, below=2.2cm of b2] (b5) {B5: use x};
\node[basic block, below=1.0cm of b5] (b6) {B6: exit};
\draw[cfg edge] (b1) -- (b2);
\draw[cfg edge] (b2) -- (b3);
\draw[cfg edge] (b2) -- (b4);
\draw[cfg edge] (b3) -- (b5);
\draw[cfg edge] (b4) -- (b5);
\draw[cfg edge] (b5) -- (b6);
\draw[back edge] (b5) -- (b2);
\end{tikzpicture}
\caption{A branch inside a loop, where \code{x} needs phis at both the branch join $B5$ and the loop head $B2$ after frontier iteration.}
\label{fig:ch12-ssa}
\end{figure}

\begin{table}
\caption{Dominance frontiers and resulting phi placement for \code{x} and \code{y} on the loop-and-branch graph.}
\label{tab:ch12-df}
\begin{tabular}{lcc}
\toprule
Block & Frontier set & Phi placed \\
\midrule
B1 & empty & none \\
B2 & loop head B2 & phi for x at B2 \\
B3 & join B5 & phi for x at B5 \\
B4 & join B5 & phi already at B5 \\
B5 & loop head B2 & none new \\
\bottomrule
\end{tabular}
\end{table}

Placement for \code{x} seeds the worklist with $B1$, $B3$, and $B4$.
Processing $B1$ adds nothing because its frontier is empty.
Processing $B3$ adds a phi at $B5$ and queues $B5$.
Processing $B4$ finds $B5$ already placed and adds nothing.
Processing $B5$ adds a phi at $B2$ and queues $B2$.
Processing the re-seeded $B2$ finds its own frontier already satisfied.
The final placement holds a two-argument phi at $B5$ for the branch and a two-argument phi at $B2$ for loop entry versus back edge.
Placement for \code{y} seeds with $B3$ alone, finds frontier entries $B5$ and $B2$, and would place phis there in minimal form.
Pruning removes both \code{y} phis because no use of \code{y} is live across either join in this example.
The surviving program therefore merges only \code{x}, which matches the intuition that unmerged variables need no phis.

Renaming starts at $B1$ with empty stacks and pushes $x1$ for the definition there.
It descends to $B2$, creates the head phi result $x2$, and pushes it for the whole loop body.
Along the left edge it pushes $x3$ at $B3$, and along the right edge it pushes $x4$ at $B4$.
At $B5$ it creates $x5$ as $\phifn{x_3, x_4}$ using the edge-specific tops, then rewrites the use to $x5$.
The back edge carries $x5$ to the head phi, whose other argument reads $x1$ from entry.
Because $B6$ is a dominator-tree child of $B5$, it is visited while $x5$ is still on the stack.
A use of \code{x} at the exit therefore reads $x5$, the value leaving the final iteration, not the pre-loop value $x1$.
Figure~\ref{fig:visual-ssa-domtree} shows why: $B6$ is below $B5$ in the dominator tree, so the renaming traversal has not yet popped $B5$'s version.
Compare this tree with the execution edges in Figure~\ref{fig:ch12-ssa}; the two graphs answer different questions.

\begin{figure}[htbp]
  \centering
  \begin{tikzpicture}[node distance=11mm and 10mm]
    \node[basic block] (b1) {$B1$: define $x_1$};
    \node[basic block,below=of b1] (b2) {$B2$: define $x_2$};
    \node[basic block,below=of b2] (b4) {$B4$: define $x_4$};
    \node[basic block,left=of b4] (b3) {$B3$: define $x_3$};
    \node[basic block,right=of b4] (b5) {$B5$: define $x_5$};
    \node[basic block,below=of b5] (b6) {$B6$: use $x_5$};
    \draw[cfg edge] (b1) -- (b2);
    \draw[cfg edge] (b2) -- (b3);
    \draw[cfg edge] (b2) -- (b4);
    \draw[cfg edge] (b2) -- (b5);
    \draw[cfg edge] (b5) -- (b6);
    \node[font=\sffamily\small,align=left,below=15mm of b3,anchor=north]
      {Edges are immediate-dominator links,\\not execution transfers.};
  \end{tikzpicture}
  \caption{The dominator-tree walk visits the exit while the join's version $x_5$ is still active; it does not restore the pre-loop version first.}
  \label{fig:visual-ssa-domtree}
\end{figure}


The \enquote{one definition per name} invariant holds at every step, and each use names the version from its own path.
A checker can verify the result linearly by confirming that each version has one definition and each phi has one argument per incoming edge.

\section{Implementation notes and cost}
\label{sec:ch12-implementation}

A builder computes dominators, derives frontiers, places phis per variable, and renames along the dominator tree.
For $N$ blocks and $E$ edges, the simple upward-walk frontier construction takes at most $O(EN)$ time with constant-time set insertion.
If $F$ is the sum of frontier sizes and $V$ the number of variables, placement takes $O(V(N+F))$ time with duplicate-free worklists.
Renaming takes $O(N+I+P)$ time for $I$ ordinary operand occurrences and definitions and $P$ phi arguments.
Memory holds one stack entry per live variable along the current tree path, plus one phi node per placed merge.
Large functions with many globals need care: a variable defined in hundreds of blocks seeds hundreds of frontier walks.
SSA construction frameworks therefore filter with address-taken and single-block tests before placing anything.
Safonov discusses keeping such construction passes predictable inside trusted toolchains, see~\cite{safonov-2010-trustworthy}.

Destruction belongs in the notes because construction without a destruction plan strands phis before emission.
Critical edges, where a block with several successors jumps to a block with several predecessors, need splitting so phi copies land on exactly one edge.
Lost-copy and swap problems arise when parallel phi semantics meet sequential moves, and the standard fix orders copies through temporaries along each edge.
Register allocation either coalesces phi-related copies aggressively or lowers phis before coloring, depending on the allocator architecture.
The C standard constrains observable behavior across these rewrites, so copies that touch volatile locations or sequence points need conservative handling, see~\cite{wg14-2011-c11}.
Cytron and colleagues develop dominance-frontier-based SSA construction and its connection to control dependence; read the construction and its proof together~\cite{cytron-1991-ssa}.

\section{Exercises}
\label{sec:ch12-exercises}

\begin{enumerate}
\item \label{ex:ch12-place-rename} Place phis and rename Figure~\ref{fig:ch12-ssa} by hand after adding a definition of \code{y} in $B4$.
Use Table~\ref{tab:ch12-df}; adding definitions without changing edges does not change dominance frontiers.
Show the version stacks on entry to every block and the final phi arguments at $B2$ and $B5$.
State which \code{y} phis pruning keeps when $B6$ uses \code{y}.
\item \label{ex:ch12-critical-edge} Find a critical edge in Figure~\ref{fig:ch12-ssa} after adding an edge from $B3$ directly to $B6$.
Explain why phi copies for $B5$ cannot sit in $B3$ or $B5$ without affecting the new edge.
Split the critical edge, show where the copies land, and run Algorithm~\ref{alg:ch12-placement} again to confirm the placement is unchanged.
\item \label{ex:ch12-lost-copy} Construct a two-predecessor join where naive sequential phi lowering swaps two variables instead of copying them.
Show the wrong sequential moves, the correct parallel semantics with $\assignphi{}{}$ notation, and the temporary-based fix.
Implement the fix for arbitrary phi arity and test it on random joins with three predecessors.
\end{enumerate}

\section{Further reading}
\label{sec:ch12-reading}

Appel gives the stack-based renaming walk in implementable detail, see~\cite{appel-1998-modern}.
Morgan places SSA construction among the loop and redundancy passes that consume it, see~\cite{morgan-1998-optimizing}.
Safonov covers the engineering discipline that keeps construction and destruction predictable, see~\cite{safonov-2010-trustworthy}.
The C language definition governs which copies and merges an implementation may reorder or remove, see~\cite{wg14-2011-c11}.
